% \usepackage{footnote}
\usepackage{float}
\usepackage{url}

% High quality fonts
\usepackage[T1]{fontenc}

% Underlining text
\usepackage[normalem]{ulem}

% Math options
\usepackage{amsmath}
\usepackage{amsthm}
\usepackage{nicefrac}

% Theorem naming
\newtheorem{lem}{Lemma}
\newtheorem{proposition}{Proposition}

% Special tikz drawings
\usepackage{tikz}
\def\checkmark{\tikz\fill[scale=0.4](0,.35) -- (.25,0) -- (1,.7) -- (.25,.15) -- cycle;}

% Caption options
\usepackage{caption}
\captionsetup{labelfont=bf, margin=1cm, tableposition=top, figurewithin=none, tablewithin=none}

\raggedbottom

% For tables
\usepackage{booktabs}
\usepackage{siunitx}
\usepackage{makecell}

% For graphs
\usepackage[all,cmtip]{xy}

\usepackage{xstring}
\usepackage{ifthen}
\usepackage{catchfile}

\newcommand{\scalarinput}[1]{$\input{../output/tables/scalars/#1}$} % Keep minus signs
\newcommand{\absinput}[1]{%
    \CatchFileDef{\filecontents}{../output/tables/scalars/#1}{}% Read the file contents into \filecontents
    \IfBeginWith{\filecontents}{-}{% Check if the first character is a minus sign
        \StrGobbleLeft{\filecontents}{1}[\filecontents]% Remove the first character
    }{}%
    $\filecontents$% Output the processed contents
}
\newcommand{\scalarpctinput}[1]{%
  \CatchFileDef{\filecontents}{../output/tables/scalars/#1}{}% Read the file contents into \filecontents
  \IfBeginWith{\filecontents}{0.}{% Check if the first character is a minus sign
      \StrGobbleLeft{\filecontents}{2}[\filecontents]% Remove the first character
  }{}%
  $\filecontents$% Output the processed contents
}
% For xmark
\usepackage{pifont}% http://ctan.org/pkg/pifont
\newcommand{\xmark}{\text{\ding{55}}}

% Columns
\usepackage{multicol}

% No figures below footnote
\usepackage[bottom]{footmisc}

% Adding notes to table
\newcommand{\tablenote}[1]  {\par\vspace{3pt}{\raggedright \scriptsize \emph{Notes: }#1\par}\vspace{2pt}}

% Colors
% These are my colors -- there are many like them, but these ones are mine.
\definecolor{blue}{RGB}{0,92,216}
\definecolor{red}{RGB}{213,40,0}
\definecolor{yellow}{RGB}{240,228,66}
\definecolor{green}{RGB}{0,158,115}
\definecolor{grey}{RGB}{180,180,180}
\definecolor{darkgrey}{RGB}{105,105,105}
\definecolor{white}{RGB}{255,255,255}
\definecolor{purple}{RGB}{199,124,255}
\definecolor{orange}{RGB}{197,90,17}

\newcommand{\grey}[1]{\textcolor{grey}{#1}}
\newcommand{\red}[1]{\textcolor{red}{#1}}
\newcommand{\green}[1]{\textcolor{green}{#1}}
\newcommand{\blue}[1]{\textcolor{blue}{#1}}
\newcommand{\purple}[1]{\textcolor{purple}{#1}}
\newcommand{\orange}[1]{\textcolor{orange}{#1}}

% Useful math commands
\newcommand{\argmax}[0]{\operatornamewithlimits{argmax}}
\newcommand{\argmin}[0]{\operatornamewithlimits{argmin}}
\newcommand{\E}{\mathbb{E}}
\newcommand{\R}{\mathbb{R}}
\renewcommand{\d}{\mbox{ d}}
\newcommand{\Var}{\mbox{Var}}

% Font -- seems to need to be after things
\usepackage[sc]{mathpazo}

% Bibliography
\usepackage[authordate, backend=biber, maxcitenames=2, mincitenames=1, maxbibnames=99, uniquelist=false, uniquename=false]{biblatex-chicago}
\makeatletter
\renewcommand{\@makefntext}[1]{%
  \setlength{\parindent}{0pt}%
  \makebox[1.5em][r]{\textsuperscript{\@thefnmark}}#1}
\makeatother

% Remove extraneous fields
\AtEveryBibitem{%
  \clearfield{doi}%
  \clearfield{issn}%
  \clearfield{isbn}%
  \clearfield{url}%
  \clearfield{urldate}%
  \clearfield{eprint}%
  \clearfield{note}%
}

% Subfiles -- needs to be near end
\usepackage{subfiles}

% For crossing things out
\usepackage{centernot}
\newcommand{\notimplies}{\centernot\implies}

% Math commands
\renewcommand{\lvert}[1]{\left.#1\right|}
\renewcommand{\rvert}[1]{\left|#1\right.}

% For separating out referee comments
\newenvironment{refcomment}
  {\clearpage\begin{quotation}\color{blue}\singlespacing\small}
{\end{quotation}}

\newenvironment{refcommentnoclear}
  {\begin{quotation}\color{blue}\singlespacing\small}
{\end{quotation}}

% For marking LLM-generated text (displays in red for review)
\newenvironment{llm}
  {\color{red}}
  {}

% For side by side graphs
\usepackage{subcaption}

% Coloring tables
\usepackage{colortbl}

% For inputting specific things about the counterfactuals
\newcommand{\passval}[1]{\scalarinput{#1_1}}
\newcommand{\passabs}[1]{\absinput{#1_1}}
\newcommand{\passse}[1]{(\scalarinput{#1_1_se})}
\newcommand{\passsetext}[1]{\scalarinput{#1_1_se}}

\newcommand{\nopassval}[1]{\scalarinput{#1_0}}
\newcommand{\nopassabs}[1]{\absinput{#1_0}}
\newcommand{\nopassse}[1]{(\scalarinput{#1_0_se})}
\newcommand{\nopasssetext}[1]{\scalarinput{#1_0_se}}
